\chapter{Métaordres, liquidité latente et impact croisé}\label{ch:mx:metaorders-latent-liquidity-and-cross-impact}

Pourquoi l'impact devrait-il croître comme la racine carrée de la taille, quel que soit le titre, la décennie ou le marché~? La réponse que donnent aujourd'hui la plupart des modèles est que le carnet d'ordres
que voit un trader ne représente qu'une petite part de la liquidité existante~: le reste est latent, en attente d'être révélé à des prix que n'a pas encore atteints un ordre de grande taille. Ce
chapitre construit ce carnet latent, retrouve la racine carrée et le régime linéaire situé en dessous, vérifie à quelles conditions un métaordre est tarifé équitablement, et aborde le
impact qu'a la négociation d'un titre sur un autre.

\section{Pourquoi l'impact est concave}

Le chapitre~11 a mesuré un impact concave, et le chapitre~12 un impact transitoire. Trois explications de la concavité sont utilisées. La fragmentation~: les grands métaordres sont découpés
en fils dont les impacts se dissipent entre eux. L'information~: Farmer, Gerig, Lillo et Waelbroeck (2013) ont dérivé la forme de l'impact à partir de la concurrence entre
services d'exécution et d'une \omterm{def:mx:metaorders-latent-liquidity-and-cross-impact:fair}{condition de prix équitable}. La liquidité~: Toth et ses coauteurs (2011, chapitre~11) ont soutenu que l'offre et la demande moyennes s'annulent au
prix et croissent linéairement en s'en éloignant. La troisième explication est devenue un modèle.

\section{Liquidité latente et carnet d'ordres localement linéaire}

\begin{definition}[Liquidité latente, carnet d'ordres latent]\label{def:mx:metaorders-latent-liquidity-and-cross-impact:latent}
La \emph{liquidité latente}\index{liquidité latente} est l'intention de négocier à un prix qui n'a pas été placé dans le carnet, parce que son détenteur attend
que le prix s'en approche. Le \emph{carnet d'ordres latent}\index{carnet d'ordres latent} est la densité de ces intentions par prix, dont le carnet visible
constitue la petite partie la plus proche du prix.
\end{definition}

\begin{definition}[Carnet d'ordres localement linéaire]\label{def:mx:metaorders-latent-liquidity-and-cross-impact:linear}
Un \emph{carnet d'ordres localement linéaire}\index{carnet d'ordres localement linéaire} possède une densité latente qui s'annule au prix et croît linéairement en s'en éloignant,
$\varphi(x)=L\,|x-p|$, où $L$ désigne la liquidité par unité de prix au carré.
\end{definition}

Donier, Bonart, Mastromatteo et Bouchaud (2015) l'ont dérivé de la réaction--diffusion des ordres latents~: les intentions des acheteurs et des vendeurs diffusent au fil
de l'évolution de leurs valorisations, et s'annihilent (se traitent) là où elles se rencontrent. La densité signée $\varphi$ (les cours vendeurs étant positifs) obéit alors à une équation de diffusion
dont la solution stationnaire est linéaire. Un métaordre d'achat au taux $m$ retire des cours vendeurs au prix, et le prix, zéro de $\varphi$, monte. Si le métaordre est
rapide par rapport à la diffusion, il consomme simplement le carnet linéaire~: déplacer le prix de $I$ consomme $\int_0^IL\,x\,dx=LI^2/2$, donc $I=\sqrt{2Q/L}$. S'il est
lent, le carnet se reconstitue par diffusion au fur et à mesure de la négociation, et l'impact est linéaire en taille. Bucci, Benzaquen, Lillo et Bouchaud (2019) ont documenté ce passage d'un régime à l'autre.

\texttt{firm.crossimpact.llob} intègre la diffusion sur une grille et retire, à chaque pas, les actions du métaordre parmi les cours vendeurs les plus proches du prix.

\omcode{code/firm/crossimpact/firm_crossimpact.py}{33}{56}{Le carnet d'ordres latent~: la densité signée diffuse, le métaordre retire les cours vendeurs les plus proches, et le prix est le zéro de la densité.}\label{lst:mx:metaorders-latent-liquidity-and-cross-impact:llob}

Avec $L=1$, $D=1$ et une exécution de dix secondes, l'\omterm{def:mx:the-empirics-of-price-impact:peak}{impact maximal} vaut 0.05 tick pour $Q=0.3$, 0.18 pour 1, 0.52 pour 3, 1.74 pour 10, 4.9 pour 30, 12.2 pour 100, 23.3 pour 300 et
44.0 pour 1\,000 (\cref{fig:mx:metaorders-latent-liquidity-and-cross-impact:scan}). L'exposant local vaut 1.0 jusqu'à $Q=10$, puis 0.94, 0.76, 0.59 et 0.53~: linéaire
en dessous de $LDT=10$, racine carrée au-dessus, où l'impact tend vers $\sqrt{2Q/L}$ (24.5 pour 300, 44.7 pour 1\,000).

\begin{omfigure}
\begin{tikzpicture}
\begin{axis}[width=9cm,height=5.4cm,xmode=log,ymode=log,xlabel={\small taille du métaordre $Q$},ylabel={\small impact maximal (ticks)},
  tick label style={font=\footnotesize},grid=major,grid style={gray!25},legend columns=3,legend style={font=\footnotesize,at={(0.5,-0.3)},
  anchor=north,draw=none,fill=none,/tikz/every even column/.append style={column sep=8pt}},table/col sep=comma]
\addplot[omDef,very thick,mark=*] table[x=q,y=peak]{figdata/microstructure/13-metaorders-latent-liquidity-and-cross-impact/scan.csv};
\addplot[gray,dashed,domain=0.3:30] {0.176*x};
\addplot[omThm,dashed,domain=30:1000] {sqrt(2*x)};
\legend{carnet latent,linéaire,$\sqrt{2Q/L}$}
\end{axis}
\end{tikzpicture}

\omcaption{\omterm{def:mx:the-empirics-of-price-impact:peak}{Impact maximal} d'un métaordre de dix secondes sur un carnet latent localement linéaire ($L=1$, $D=1$)~: linéaire en taille lorsque le carnet se reconstitue plus vite que
l'ordre ne le consomme, racine carrée dans le cas contraire. Données~: \texttt{mx\_cross.llob\_scan}.}\label{fig:mx:metaorders-latent-liquidity-and-cross-impact:scan}
\end{omfigure}

\section{Tarification équitable et fin d'un métaordre}

\begin{definition}[Condition de prix équitable]\label{def:mx:metaorders-latent-liquidity-and-cross-impact:fair}
La \emph{condition de prix équitable}\index{condition de prix équitable} (fair pricing) (Farmer, Gerig, Lillo et Waelbroeck, 2013) énonce que le prix moyen payé pour un
métaordre est égal au prix observé après son achèvement~: celui qui l'a exécuté n'a pas été rémunéré pour s'être fait prélever, et le détenteur du métaordre a payé ce que sa
négociation a révélé.
\end{definition}

Si l'impact croît comme la racine carrée du temps de négociation, l'impact moyen payé pendant l'exécution vaut les deux tiers du pic, de sorte que la \omterm{def:mx:metaorders-latent-liquidity-and-cross-impact:fair}{condition de prix équitable} prédit
que le prix se stabilise aux deux tiers de son pic une fois le métaordre terminé. Sur le carnet latent, un métaordre de 10 atteint 1.74 tick et coûte en moyenne 1.16,
soit 0.66 du pic (\cref{fig:mx:metaorders-latent-liquidity-and-cross-impact:path})~; après la fin, le prix repasse sous la moyenne payée 1.9 seconde
plus tard, et continue de baisser, jusqu'à 0.16 du pic après 100 secondes. Pour un métaordre de 300, la moyenne vaut de nouveau 0.66 du pic, franchie 2.3 secondes après
la fin, avec 0.27 restant après 100 secondes. Le modèle sans perte de mémoire rend l'impact transitoire, si bien que le prix n'est équitable que l'espace d'un instant~; dans le modèle complet
une mémoire finie des ordres latents laisse une part permanente.

\begin{omfigure}
\begin{tikzpicture}
\begin{axis}[width=9cm,height=5cm,xlabel={\small secondes},ylabel={\small impact (en ticks)},tick label style={font=\footnotesize},xmin=0,xmax=110,
  ymin=0,ymax=2,grid=major,grid style={gray!25},table/col sep=comma]
\fill[gray!12] (axis cs:0,0) rectangle (axis cs:10,2);
\addplot[omDef,very thick] table[x=t,y=p]{figdata/microstructure/13-metaorders-latent-liquidity-and-cross-impact/path.csv};
\addplot[omThm,dashed,domain=0:110] {1.155};
\node[omThm,font=\footnotesize,anchor=south west] at (axis cs:30,1.16) {prix moyen payé};
\end{axis}
\end{tikzpicture}

\omcaption{Un métaordre de 10 sur le carnet latent~: l'impact pendant l'exécution (en grisé) et après celle-ci, comparé au prix moyen payé. Le prix
franchit la moyenne deux secondes après la fin et continue de baisser. Données~: \texttt{mx\_cross.fair}.}\label{fig:mx:metaorders-latent-liquidity-and-cross-impact:path}
\end{omfigure}

\section{Impact croisé}

\begin{definition}[Impact croisé, matrice d'impact croisé]\label{def:mx:metaorders-latent-liquidity-and-cross-impact:cross}
L'\emph{impact croisé}\index{impact croisé} est l'effet de la négociation d'un actif sur le prix d'un autre. La \emph{matrice d'impact croisé}\index{matrice d'impact croisé}
$\Lambda$ donne la variation de prix de chaque actif par unité du flux d'ordres de chaque actif, $r_t=\Lambda q_t+\text{bruit}$.
\end{definition}

Pasquariello et Vega (2015) ont constaté sur les actions américaines que le déséquilibre quotidien des ordres d'une branche ou d'une action déplace les rendements d'autres branches et d'autres actions,
de manière persistante et souvent négative. Benzaquen, Mastromatteo, Eisler et Bouchaud (2017) ont construit un propagateur multivarié dans lequel la liquidité partage la
structure sectorielle des corrélations, et ont montré que les transactions médiatisent une part significative de la covariance des rendements.

Deux actions simulées dont les flux d'ordres sont corrélés à 0.6 (en milliers d'actions par intervalle) et une matrice vraie $\Lambda=\begin{pmatrix}1.0&0.4\\0.4&0.8\end{pmatrix}$
ticks par millier d'actions, avec un bruit unitaire, donnent sur 2\,000 intervalles l'estimation $\begin{pmatrix}0.96&0.42\\0.42&0.81\end{pmatrix}$. Au-delà de son propre flux,
le flux de l'autre action explique 4.5~\% de la variance des rendements de la première et 5.3~\% de celle de la seconde~: les flux corrélés portent déjà l'essentiel du
mouvement commun, et l'\omterm{def:mx:metaorders-latent-liquidity-and-cross-impact:cross}{impact croisé} ajoute le reste.

\section{Impact croisé sans arbitrage}

Schneider et Lillo (2019) ont étendu l'argument d'absence d'arbitrage dynamique à plusieurs actifs~: avec des noyaux de décroissance bornés, l'\omterm{def:mx:metaorders-latent-liquidity-and-cross-impact:cross}{impact croisé} doit être linéaire et impair en
intensité de négociation, et symétrique, l'impact de $i$ sur $j$ étant égal à celui de $j$ sur $i$~; sinon un trader pourrait parcourir les actifs en boucle et réaliser un profit.
Une matrice estimée n'est pas exactement symétrique~; \texttt{symmetric\_psd} la projette sur la matrice symétrique semi-définie positive la plus proche avant que
tout optimiseur ne l'utilise.

Que change l'\omterm{def:mx:metaorders-latent-liquidity-and-cross-impact:cross}{impact croisé} pour l'exécution~? Vendre 10 (milliers d'actions) de chaque action sur une unité de temps, avec un impact transitoire décroissant au taux 3,
le calendrier conjoint le moins coûteux coûte 52.1 (\cref{fig:mx:metaorders-latent-liquidity-and-cross-impact:costs}). Un modèle d'impact propre, qui ignore les
termes croisés, estime ce même calendrier à 36.0, soit 31~\% de moins. Pour un trade de paire, vendre 10 de la première et acheter 10 de la seconde, l'optimum conjoint coûte
20.0 et l'estimation d'impact propre s'élève à 36.0, soit 80~\% de trop~: l'\omterm{def:mx:metaorders-latent-liquidity-and-cross-impact:cross}{impact croisé} de chaque jambe aide l'autre. Lorsque l'\omterm{def:mx:metaorders-latent-liquidity-and-cross-impact:cross}{impact croisé} décroît comme l'impact propre, les
calendriers optimaux sont les mêmes avec ou sans lui~; ce en quoi le modèle se trompe, c'est le coût. Les calendriers diffèrent lorsque les jambes sont traitées l'une après l'autre, comme
le ferait un desk qui croit les jambes indépendantes~: le calendrier conjoint permet alors d'économiser 19~\% du coût séquentiel pour la vente et 48~\% pour le trade de paire. Lorsque
l'\omterm{def:mx:metaorders-latent-liquidity-and-cross-impact:cross}{impact croisé} décroît plus lentement que l'impact propre (0.3 contre 3), l'optimum conjoint change aussi de forme et fait économiser 3.7~\% de plus sur le trade de paire.

\omcode{code/firm/crossimpact/firm_crossimpact.py}{92}{111}{Liquidation conjointe sous impact transitoire propre et croisé, ou action par action en ignorant les termes croisés.}\label{lst:mx:metaorders-latent-liquidity-and-cross-impact:liq}

\begin{omfigure}
\begin{tikzpicture}
\begin{axis}[width=8.5cm,height=5cm,ybar,bar width=12pt,xtick={0,1},xticklabels={vente des deux,trade de paire},ylabel={\small coût},
  tick label style={font=\footnotesize},ymin=0,ymax=75,grid=major,grid style={gray!25},enlarge x limits=0.4,
  legend columns=3,legend style={font=\footnotesize,at={(0.5,-0.2)},anchor=north,draw=none,fill=none,
  /tikz/every even column/.append style={column sep=8pt}},table/col sep=comma]
\addplot[omDef,fill=omDef!40] table[x=k,y=joint]{figdata/microstructure/13-metaorders-latent-liquidity-and-cross-impact/costs.csv};
\addplot[omProp,fill=omProp!40] table[x=k,y=own]{figdata/microstructure/13-metaorders-latent-liquidity-and-cross-impact/costs.csv};
\addplot[omThm,fill=omThm!40] table[x=k,y=sequential]{figdata/microstructure/13-metaorders-latent-liquidity-and-cross-impact/costs.csv};
\legend{optimum conjoint,estimation d'impact propre,jambes l'une après l'autre}
\end{axis}
\end{tikzpicture}

\omcaption{Vente de deux actions, ou trade de ces actions en paire, sous \omterm{def:mx:metaorders-latent-liquidity-and-cross-impact:cross}{impact croisé}~: le coût de l'optimum conjoint, ce qu'un modèle d'impact propre en estime,
et le coût de la négociation des jambes à tour de rôle. Données~: \texttt{mx\_cross.cross\_study}.}\label{fig:mx:metaorders-latent-liquidity-and-cross-impact:costs}
\end{omfigure}

\section{Tutoriel~: vendre l'action, déplacer sa voisine}\label{tut:mx:metaorders-latent-liquidity-and-cross-impact:tut}

\begin{tutorial}
\textbf{Objectif.} Retrouver la racine carrée à partir d'un carnet latent, vérifier le prix équitable, puis estimer et utiliser l'\omterm{def:mx:metaorders-latent-liquidity-and-cross-impact:cross}{impact croisé}.
\textbf{État final~:} \cref{fig:mx:metaorders-latent-liquidity-and-cross-impact:scan,fig:mx:metaorders-latent-liquidity-and-cross-impact:path,fig:mx:metaorders-latent-liquidity-and-cross-impact:costs}.
\begin{tutsteps}
\item \textbf{Carnet latent.} \texttt{firm\_crossimpact.llob(q, horizon, L, D, dt, width, dx, after)}~; \texttt{mx\_cross.llob\_scan()} sur huit tailles.
\item \textbf{Prix équitable.} \texttt{fair\_pricing(times, price, horizon)} et \texttt{mx\_cross.fair(q)}.
\item \textbf{\omterm{def:mx:metaorders-latent-liquidity-and-cross-impact:cross}{Impact croisé}.} \texttt{two\_assets()}, \texttt{estimate(r, q)}, \texttt{symmetric\_psd}, \texttt{explained\_by\_other}.
\item \textbf{Liquidation.} \texttt{liquidation(Lambda, rho, times, X, joint, rho\_cross)} pour une vente et un trade de paire~; dessiner avec \texttt{fig\_cross.py}.
\end{tutsteps}
\textbf{Pour aller plus loin.} Ajouter l'annulation et le dépôt d'ordres latents (mémoire finie) et chercher un impact permanent~; estimer l'\omterm{def:mx:metaorders-latent-liquidity-and-cross-impact:cross}{impact croisé} sur
trois actifs et le décomposer sur les vecteurs propres de la covariance des flux.
\end{tutorial}

\section{Construction~: liquidité latente et impact croisé}\label{bld:mx:metaorders-latent-liquidity-and-cross-impact:crossimpact}

\begin{build}
\bfield{Rôle} Le mécanisme de la loi en racine carrée et l'impact multi-actifs dont ont besoin les exécutions de portefeuille (chapitres 15 et 22) et la
teneur de marché multi-actifs du livre~11.
\bfield{Interface} \texttt{llob(q, horizon, L, D, dt, width, dx, after)}, \texttt{fair\_pricing(times, price, horizon)}, \texttt{estimate(r, q)},
\texttt{symmetric\_psd(Lambda)}, \texttt{explained\_by\_other(r, q)}, \texttt{liquidation(Lambda, rho, times, X, joint, rho\_cross)}.
\bfield{Règles} Diffusion explicite (stable pour $D\,dt/dx^2\le\tfrac12$) avec extrémités linéaires fixes~; le prix est le zéro interpolé de la densité~;
\omterm{def:mx:metaorders-latent-liquidity-and-cross-impact:cross}{impact croisé} par moindres carrés des rendements sur les flux~; projection symétrique par écrêtage des valeurs propres.
\bfield{Tests d'acceptation} \texttt{code/firm/crossimpact/tests/}~: un métaordre rapide déplace le prix d'environ $\sqrt{2Q/L}$~; une matrice implantée est retrouvée~;
la projection est symétrique et positive~; la liquidation conjointe ne coûte pas plus cher que procéder action par action.
\bfield{Pour aller plus loin} Ordres latents avec dépôt et annulation~; la décomposition en liquidités propres~; \omterm{def:mx:metaorders-latent-liquidity-and-cross-impact:cross}{impact croisé} non linéaire.
\end{build}

\begin{omsources}
\item J.~D.~Farmer, A.~Gerig, F.~Lillo et H.~Waelbroeck, ``How efficiency shapes market impact'', \emph{Quantitative Finance} 13(11), 2013.
\item J.~Donier, J.~Bonart, I.~Mastromatteo et J.-P.~Bouchaud, ``A fully consistent, minimal model for non-linear market impact'', \emph{Quantitative
Finance} 15(7), 2015.
\item P.~Pasquariello et C.~Vega, ``Strategic cross-trading in the U.S. stock market'', \emph{Review of Finance} 19(1), 2015.
\item M.~Benzaquen, I.~Mastromatteo, Z.~Eisler et J.-P.~Bouchaud, ``Dissecting cross-impact on stock markets: an empirical analysis'', \emph{Journal of
Statistical Mechanics}, 2017.
\item M.~Schneider et F.~Lillo, ``Cross-impact and no-dynamic-arbitrage'', \emph{Quantitative Finance} 19(1), 2019.
\item F.~Bucci, M.~Benzaquen, F.~Lillo et J.-P.~Bouchaud, ``Crossover from linear to square-root market impact'', \emph{Physical Review Letters} 122, 2019.
\end{omsources}

\section{Exercices}

\begin{exercise}[$\star$]\label{exo:mx:metaorders-latent-liquidity-and-cross-impact:1}
Un carnet latent localement linéaire a $L=2$ (milliers d'actions par tick au carré). De combien un métaordre rapide de 50 milliers d'actions déplace-t-il le prix~?
\end{exercise}

\begin{exercise}[$\star$]\label{exo:mx:metaorders-latent-liquidity-and-cross-impact:2}
Pourquoi l'impact est-il linéaire en taille lorsque le métaordre est lent par rapport à la diffusion du carnet latent~?
\end{exercise}

\begin{exercise}[$\star$]\label{exo:mx:metaorders-latent-liquidity-and-cross-impact:3}
Montrer que si l'impact croît comme $\sqrt t$ pendant une exécution à taux constant, l'impact moyen payé vaut les deux tiers du pic.
\end{exercise}

\begin{exercise}[$\star\star$]\label{exo:mx:metaorders-latent-liquidity-and-cross-impact:4}
La \omterm{def:mx:metaorders-latent-liquidity-and-cross-impact:cross}{matrice d'impact croisé} estimée est $\begin{pmatrix}1.0&0.5\\-0.1&-0.2\end{pmatrix}$. Qu'a-t-elle de faux, et que fait la projection~?
\end{exercise}

\begin{exercise}[$\star\star$]\label{exo:mx:metaorders-latent-liquidity-and-cross-impact:5}
Pourquoi un modèle d'impact propre sous-estime-t-il le coût de la vente de deux actions corrélées et le surestime-t-il pour un trade de paire~?
\end{exercise}

\begin{exercise}[$\star\star$]\label{exo:mx:metaorders-latent-liquidity-and-cross-impact:6}
Pourquoi l'optimum conjoint est-il le même avec et sans \omterm{def:mx:metaorders-latent-liquidity-and-cross-impact:cross}{impact croisé} lorsque les deux décroissent au même taux~?
\end{exercise}

\begin{exercise}[$\star\star\star$]\label{exo:mx:metaorders-latent-liquidity-and-cross-impact:7}
\emph{Programmation.} Relancer le balayage du carnet latent avec un horizon de 100 secondes au lieu de 10. Où le point de bascule se déplace-t-il, et pourquoi~?
\end{exercise}

\begin{exercise}[$\star\star\star$]\label{exo:mx:metaorders-latent-liquidity-and-cross-impact:8}
\emph{Trouver la faille.} «~Le prix est redescendu sous notre prix moyen d'exécution deux secondes après la fin de notre ordre, donc les teneurs de marché qui nous ont exécutés ont perdu de l'argent.~»
\end{exercise}

\section{Problème~: vendre l'action, déplacer sa voisine}

\begin{problem}[{Problème du week-end --- vendre l'action, déplacer sa voisine}]\label{pb:mx:metaorders-latent-liquidity-and-cross-impact:1}
Un gérant de portefeuille doit vendre deux actions corrélées et réaliser plus tard un trade de paire sur elles. Expliquer l'impact et planifier les exécutions.

\medskip\noindent\textbf{Partie I --- \omterm{def:mx:metaorders-latent-liquidity-and-cross-impact:latent}{Liquidité latente}.}
\begin{enumerate}
\item Donner les trois explications de l'impact concave.
\item Définir le carnet latent et le carnet localement linéaire.
\item Dériver $\sqrt{2Q/L}$ pour un métaordre rapide.
\item Donner les impacts maximaux simulés et les exposants locaux, et l'emplacement du point de bascule.
\end{enumerate}

\medskip\noindent\textbf{Partie II --- Le juste prix.}
\begin{enumerate}[resume]
\item Énoncer la \omterm{def:mx:metaorders-latent-liquidity-and-cross-impact:fair}{condition de prix équitable}.
\item Résoudre l'exercice~3.
\item Quand le prix simulé croise-t-il le prix moyen payé, et que se passe-t-il ensuite~?
\item Pourquoi l'impact est-il transitoire dans ce modèle, et qu'est-ce qui rendrait une partie de celui-ci permanente~?
\end{enumerate}

\medskip\noindent\textbf{Partie III --- \omterm{def:mx:metaorders-latent-liquidity-and-cross-impact:cross}{Impact croisé}.}
\begin{enumerate}[resume]
\item Définir la \omterm{def:mx:metaorders-latent-liquidity-and-cross-impact:cross}{matrice d'impact croisé}.
\item Que trouvent Pasquariello et Vega, puis Benzaquen et coauteurs~?
\item Donner la matrice estimée et les parts expliquées.
\item Énoncer les conditions de Schneider et Lillo et dire pourquoi elles sont nécessaires.
\end{enumerate}

\medskip\noindent\textbf{Partie IV --- Exécution et verdict.}
\begin{enumerate}[resume]
\item Donner le coût conjoint et l'estimation de l'impact propre pour la vente.
\item Et pour le trade de paire.
\item Combien coûte l'exécution des jambes l'une après l'autre~?
\item Quand l'\omterm{def:mx:metaorders-latent-liquidity-and-cross-impact:cross}{impact croisé} modifie-t-il le calendrier optimal lui-même~?
\item Énoncer le \emph{résultat nommé}~: la part du mouvement de prix de chaque actif expliquée par le flux d'ordres de l'autre actif, et le coût économisé
par une liquidation tenant compte de l'\omterm{def:mx:metaorders-latent-liquidity-and-cross-impact:cross}{impact croisé}.
\item Comment le gérant doit-il exécuter la vente~?
\item Et le trade de paire~?
\item En une phrase~: pourquoi le voisin compte-t-il~?
\end{enumerate}
\end{problem}

\section{Questions d'entretien}

\begin{interviewq}[$\star$ \iqroles{chercheur}]\label{iq:mx:metaorders-latent-liquidity-and-cross-impact:1}
Expliquer la loi en racine carrée à l'aide du \omterm{def:mx:metaorders-latent-liquidity-and-cross-impact:latent}{carnet d'ordres latent} en trois phrases.
\end{interviewq}

\begin{interviewq}[$\star\star$ \iqroles{chercheur}]\label{iq:mx:metaorders-latent-liquidity-and-cross-impact:2}
Qu'est-ce que la \omterm{def:mx:metaorders-latent-liquidity-and-cross-impact:fair}{condition de prix équitable}, et que prédit-elle pour le prix après un métaordre~?
\end{interviewq}

\begin{interviewq}[$\star\star$ \iqroles{trader}]\label{iq:mx:metaorders-latent-liquidity-and-cross-impact:3}
Vous vendez une position importante sur une action. Quelles autres positions la pénalisent, et comment le sauriez-vous~?
\end{interviewq}

\begin{interviewq}[$\star\star$ \iqroles{chercheur}]\label{iq:mx:metaorders-latent-liquidity-and-cross-impact:4}
Comment estimeriez-vous une \omterm{def:mx:metaorders-latent-liquidity-and-cross-impact:cross}{matrice d'impact croisé}, et quelles contraintes imposeriez-vous~?
\end{interviewq}

\begin{interviewq}[$\star\star$ \iqroles{développeur}]\label{iq:mx:metaorders-latent-liquidity-and-cross-impact:5}
Simulez une densité qui diffuse sur une grille à frontière mobile. Qu'est-ce qui limite le pas de temps~?
\end{interviewq}

\begin{interviewq}[$\star\star\star$ \iqroles{trader, chercheur}]\label{iq:mx:metaorders-latent-liquidity-and-cross-impact:6}
Pourquoi un trade de paire peut-il coûter moins cher que la somme de ses jambes, et comment en tireriez-vous parti à l'exécution~?
\end{interviewq}
